\subsection{Notions de base}

\Exo{3-1}

% L'exo sur les dés non truqués suppose l'indépendance ! Il faut le mettre ensuite.
%\Exo{3-2}

\Exo{3-3}

\Exo{3-13}

\Exo{3-4}

%\Exo{3-16}

%\Exo{3-5}

%\Exo{3-6}

\subsection{Probabilit\'es conditionnelles, ind\'ependance}

\Exo{3-7}

\Exo{3-8}

\Exo{3-9}

\Exo{3-10}

\Exo{3-18}

% Attention ! Ind\'ependance !!
%\Exo{3-2}

%\Exo{3-15}

%\Exo{3-17}

\Exo{3-11}

\Exo{3-12}

\Exo{3-14}

{\noindent
{\bf L\'egende :}\\

\noindent \coolexo : exercice int\'eressant.\\
\boringexo : exercice d'entra\^inement.\\
\minsyndical : exercice fondamental.\\
\mortelexo : exercice d\'elicat.}
